\documentclass[11pt]{article}
\usepackage[margin=1in]{geometry}
\usepackage{amsmath,amssymb}
\usepackage{booktabs}
\usepackage{longtable}
\usepackage{array}
\usepackage{enumitem}
\usepackage[hidelinks]{hyperref}
\usepackage{xurl}
\usepackage{xcolor}
\setlist{nosep}

\newcommand{\todo}[1]{\textcolor{red}{[#1]}}

\title{Data notes: data centers in the US production and investment networks\\
\large P\&P project (spin-off of RbcProdNet)}
\author{}
\date{September 2026}

\begin{document}
\maketitle

\section*{Scope and status (read first)}
\begin{itemize}
  \item \textbf{Question.} Can we put data centers, and their energy use, into the 37-sector (41 after the energy split of \texttt{data\_notes.tex}) production and investment networks, and use the node to study the 2023-- investment boom? This note collects what the official statistics identify, what outside sources add, and how the node would be built. It does not decide how the boom is modelled; Section~\ref{sec:boom} lists the options and what the literature has done.
  \item \textbf{Two objects, not one.} A data center is (i) a \emph{capital good}: a structure plus servers/GPUs, networking, electrical and cooling equipment, owned by whoever operates it; and (ii) a \emph{producing establishment} selling computing, storage and hosting services. The investment boom is about (i); the energy use is about (ii), because electricity is an operating input of the establishment, not of the building. The network needs both: a data center \emph{node} whose capital is the data center \emph{bundle} (a column of the investment network) and whose main intermediate input is electricity (a column of the IO table).
  \item \textbf{Where it is in the data: not in Real estate.} Data center establishments are NAICS 518210 (colocation, cloud infrastructure, hosting; the 2022 NAICS manual sends colocation explicitly to 518210, not to lessors 531120). They sit in BEA industry 514 ``Data processing, internet publishing, and other information services'', i.e.\ inside our sector 26 \texttt{Info}. Hyperscaler data centers are own-account capital of their parents' industries (Amazon/AWS, Alphabet, Meta in 514; Microsoft and Oracle in 511 software publishers); colocation landlords Equinix and Digital Realty are REITs for tax and SEC purposes but 518210 by activity; enterprise data centers are own-account capital of every industry. The split is therefore of \texttt{Info}, with a secondary adjustment to \texttt{Real\_estate}, whose landlords buy tenants' electricity (Section~\ref{sec:where}).
  \item \textbf{What BEA now publishes.} Since the 25 September 2025 annual update, NIPA table family 5.4.x has a ``Data centers'' line under office structures (line 5 of 5.4.5U, series LA001282), quarterly from 2020Q1, sourced from the Census construction-spending data center category (published since May 2024, monthly back to January 2014). This is the \emph{structures} part only: \$20bn in 2023, \$35bn in 2024, \$50bn in 2025 (Census annual), about \$70bn at annual rate by mid-2026. There is no data center asset type in the fixed-asset accounts by industry (last update 26 Sept 2025), no data center industry in the IO accounts, and no data center equipment series. The Fed (Brandsaas et al.\ 2025) imputes equipment as computers investment above its 2015--2022 trend and finds structures $\approx 30\%$, equipment $\approx 70\%$, with 80\% of the equipment imported (Section~\ref{sec:fa}).
  \item \textbf{Energy.} Data centers run on grid electricity: 176~TWh in 2023 (4.4\% of US electricity), 192~TWh in 2024 (4.7\%), projected 325--580~TWh in 2028 and 520--840~TWh in 2030 (LBNL). On-site gas generation is about 2~GW online in mid-2026 against roughly 50~GW of data center load, with 80~GW announced; nuclear and renewable PPAs are contracts, not physical supply. In an IO table the node buys from the electric power sector, and the fuel behind that electricity is the generation mix of the grids where data centers sit (Virginia and Ohio about 60\% gas, against 43\% for the US). Electricity is about 7\% of the all-in annual cost of an AI data center, 25--30\% of the cost of the facility excluding IT, 15--27\% of colocation revenue, and about 1\% of the gross output of industry 514 (Sections~\ref{sec:energy}--\ref{sec:cost}).
  \item \textbf{Verdict.} Feasible in two months as a \emph{single} data center node carved out of \texttt{Info}, defined by function (all US data center capacity regardless of owner) with a 2023 or 2024 snapshot steady state: gross output from capacity times engineering cost per MW, IO column from the 2017 detailed Use table for 518200 re-weighted to an engineering cost structure, investment column from BEA fixed assets for industry 5140 plus the NIPA data center structures line, sales row to \texttt{Info}, \texttt{Prof\_Tech}, \texttt{Finance}, \texttt{Retail} and consumption. Not feasible officially: a time series of data center output or TFP (industry 514 is payment processing and internet publishing before 2015), and a separation of hyperscaler own-use capacity from their parents' other capital. Two modelling issues must be settled before the data work: how the boom is triggered (news about future productivity is the closest fit to the timing) and how to treat imported IT equipment in a closed-economy model in which sector 9 \texttt{Computers} would otherwise produce it (Section~\ref{sec:design}). The current answer to the second is ``US+Nvidia'': sector 9 augmented with Nvidia's US-installed data-center revenue as a factoryless producer, with the foundry, memory and assembly content as an imported non-produced input (Section~\ref{sec:nowcast} and \texttt{data\_center\_nowcast\_plan.md}).
\end{itemize}

\section{Where data centers are in the official statistics}\label{sec:where}

\subsection{Classification}
\begin{itemize}
  \item \textbf{NAICS 2017 518210} ``Data processing, hosting, and related services'': web and application hosting, data storage, and data processing (payroll, payment, and record processing). \textbf{NAICS 2022 518210} ``Computing infrastructure providers, data processing, web hosting, and related services'' adds IaaS/PaaS, cloud computing and storage, colocation, and cryptocurrency mining explicitly; it moves on-site facilities management to 541513 and payroll to 541214; internet publishing (519130) goes to 513 and 516210 and web search to 519290. The 2022 manual entry for 531120 (lessors of nonresidential buildings) states that colocation is classified in 518210. Economic Census product lines for 518210 include ``colocation services'' and ``cloud computing and cloud storage''. \url{https://www.census.gov/naics/reference_files_tools/2022_NAICS_Manual.pdf}
  \item \textbf{BEA industry codes.} Detailed benchmark (about 400 industries, 2007/2012/2017): 518200 data processing and hosting; 519130 internet publishing and web search; 511200 software publishers; 531ORE other real estate. Summary (71 industries, annual 1997--2023): 514 $=$ 518 $+$ 519; 511; ORE. Our sector 26 \texttt{Info} $=$ 511 $+$ 512 $+$ 513 $+$ 514. BEA's industry accounts remain on NAICS 2017 through the 2023 comprehensive update (the 2022 NAICS 516 split is a stated conversion problem); the 2023 update also moved equity REITs from finance into real estate.
  \item \textbf{Company vs establishment.} SEC SIC codes: Equinix and Digital Realty 6798 (REIT); Amazon 5961 (retail); Microsoft, Oracle, CoreWeave 7372 (software); Alphabet, Meta 7370. BEA's fixed assets by industry are on a company basis, so own-account servers bought by Microsoft are investment of 5110 and by Alphabet, Meta and Amazon-AWS of 5140 (Amazon's fulfilment capital is 4540/retail). The IO accounts are establishment-based, so a hyperscaler's data center establishment is 518210 whatever the parent's SIC. Enterprise data centers never appear as data centers anywhere.
  \item \textbf{Ownership of US capacity.} FERC's 2025 State of the Markets puts about 50~GW of data center capacity in service at end-2025 (24\% annual growth since 2020). CBRE counts 9.4~GW of wholesale colocation inventory in eight primary markets at end-2025 (Northern Virginia 4.0~GW) and 6.0~GW under construction, excluding hyperscaler-owned campuses. Synergy's global split at end-2025: hyperscale 48\% of capacity (60\% owned, 40\% leased), other colocation 20\%, enterprise on-premises 32\%; the US is more hyperscale-owned than that average. Applied to 50~GW: hyperscale-owned $\approx 14$~GW, hyperscale-leased $\approx 10$, other colocation $\approx 10$, enterprise $\approx 16$ (illustrative). LBNL: 74\% of servers were in hyperscale and colocation facilities in 2023, 85\% by 2028, and more than 90\% of server energy.
\end{itemize}

\subsection{Input-output accounts: what industry 518200 buys and sells}\label{sec:io}
Extracted from \texttt{IOUse\_After\_Redefinitions\_PRO\_2017\_Detail.xlsx} and the 71-industry annual Use tables (BEA \texttt{AllTablesIO.zip}, \url{https://apps.bea.gov/industry/iTables%20Static%20Files/AllTablesIO.zip}); producers' prices, \$ millions.

\begin{center}\small
\begin{tabular}{lrrr}
\toprule
Industry 518200, detailed Use table & 2007 & 2012 & 2017 \\
\midrule
Gross output (intermediates $+$ value added) & 87{,}626 & 117{,}656 & 224{,}377 \\
Total intermediate inputs & 34{,}142 & 58{,}121 & 91{,}991 \\
Compensation of employees & 25{,}671 & 18{,}803 & 43{,}130 \\
Gross operating surplus & 26{,}493 & 39{,}104 & 86{,}201 \\
\addlinespace
Purchases of electric power (221100) & 556 & 881 & 2{,}215 \\
\quad as \% of gross output & 0.63 & 0.75 & 0.99 \\
Purchases of natural gas distribution (221200) & -- & -- & 5 \\
Purchases of other real estate (531ORE) & 1{,}799 & 3{,}050 & 7{,}541 \\
Purchases of wired telecom (517110) & 1{,}250 & 1{,}256 & 1{,}284 \\
Purchases of software publishers (511200) & -- & -- & 1{,}220 \\
Purchases of computer systems design (541512) & -- & -- & 4{,}060 \\
\bottomrule
\end{tabular}
\end{center}

Largest intermediate inputs of 518200 in 2017: other real estate (7.5bn), advertising (7.0), employment services (6.6), broadcasting (6.1), lessors of intangibles (5.3), management consulting (4.7), computer systems design (4.1); electricity ranks 18th. Value added is 59\% of gross output and gross operating surplus 38\%, consistent with a capital-intensive industry; but an electricity share of 1\% says that in 2017 the industry's output was dominated by activities other than running servers. Census SAS confirms it: purchased electricity of NAICS 518 was \$980m in 2017 (0.8\% of \$127bn of expenses; the series ends in 2017). For comparison, US data center electricity use in 2017 was about 100~TWh (\$7--8bn at commercial rates), of which 518200 records 2.2bn, 519130 (Google, Meta) 0.7bn and 511200 (Microsoft) 0.6bn; the rest sits in 531ORE, which buys \$36bn of electricity (13\% of all intermediate electricity use) because landlords pay tenants' power, and in enterprise data centers in every industry.

\textbf{Make table 2017.} Commodity 518200 (total 248bn) is produced by industry 518200 (215bn) and, as a secondary product, by software publishers 511200 (33bn: Microsoft's cloud); real estate produces none of it, so BEA does not treat colocation as real estate output. Industry 518200 also produces 7.8bn of software.

\textbf{Who buys 518200 output (2017 row):} intermediate use 171bn, personal consumption 32bn. Top buyers: nondepository credit intermediation (23.9bn, 14\% of intermediate use: card processing), state and local government (14.2 $+$ 7.8bn), employment services (11.4), depository credit (9.1), travel reservation (4.5), restaurants (4.4), electric power (4.4), motor vehicle dealers (4.0). This is the row of a business-services industry; a data center node built from it inherits payment-processing customers unless the row is re-weighted with cloud revenue by customer industry.

\begin{center}\footnotesize
\begin{tabular}{lrrrrrrr}
\toprule
Industry 514, summary Use table (\$bn) & 1997 & 2007 & 2012 & 2017 & 2020 & 2022 & 2023 \\
\midrule
Gross output & 47.3 & 126.1 & 211.0 & 400.2 & 551.4 & 750.0 & 802.1 \\
Intermediate inputs & 13.4 & 44.5 & 99.2 & 159.6 & 229.3 & 333.0 & 319.2 \\
Compensation & 27.6 & 38.3 & 38.1 & 91.7 & 108.1 & 136.2 & 156.6 \\
Gross operating surplus & 5.7 & 41.1 & 71.1 & 143.4 & 209.5 & 272.1 & 317.3 \\
Purchases of utilities (22) & 0.18 & 0.60 & 1.22 & 3.04 & 5.65 & 10.33 & 8.26 \\
\quad as \% of gross output & 0.4 & 0.5 & 0.6 & 0.8 & 1.0 & 1.4 & 1.0 \\
Sales to PCE & 2.8 & 15.9 & 28.8 & 46.7 & 93.0 & 125.8 & 141.4 \\
\midrule
Other real estate (ORE): utilities purchases & 29.0 & 54.4 & 55.4 & 42.4 & 24.6 & 40.1 & 41.5 \\
\bottomrule
\end{tabular}
\end{center}
Industry 514 grew 17-fold in nominal gross output over 1997--2023 and its utilities bill 45-fold, but the utilities share never exceeds 1.4\%. The 2023 figure (8.3bn) is 40--60\% of what LBNL's 176~TWh implies (\$14bn at industrial, \$22bn at commercial rates). Census SAS revenue of NAICS 518 alone: \$220bn (2020), \$271bn (2021), \$308bn (2022); QSS about \$329bn (2023) and \$375bn (2024). BEA's Digital Economy Satellite Account (discontinued; last release December 2023, years 2017--2022) measured ``cloud services'' gross output at \$187bn in 2021 and ``internet and data services'' at \$213bn, with cloud value added up 232\% over 2017--2022; the SAS ``hosting and IT infrastructure provisioning'' product line was \$110bn in 2017 and \$156bn in 2020.

\subsection{Fixed assets: NIPA data center structures and investment by industry}\label{sec:fa}
\begin{itemize}
  \item \textbf{NIPA 5.4.5U/5.4.6U}, line 5 ``Data centers'' (series LA001282 current \$, LB001282 chained, IA001282 price), under office within commercial and health care structures, quarterly from 2020Q1, introduced 25 Sept 2025. Source: Census Value Put in Place. Deflator: BEA composite with equal weights on the BLS PPIs for industrial buildings and warehouses (the rest of office keeps the office-building PPI). FRED still shows office as one line (\$100.5bn 2024, \$102.6bn 2025). \url{https://www.bea.gov/information-updates-national-economic-accounts-2025}; series register \url{https://www.bea.gov/sites/default/files/2025-08/NEA-AU25-series-register.txt}.
  \item \textbf{Census C30}, private nonresidential, office: data center; first published with the May 2024 release, monthly back to January 2014; definition ``buildings that contain the hardware needed for storing, processing, and transmitting digital information''; racks and servers excluded. Annual value put in place, \$bn (\url{https://www.census.gov/construction/c30/pdf/region.pdf}): 2019 8.5; 2020 9.2; 2021 9.9; 2022 12.6; 2023 20.0; 2024 34.8; 2025 49.7 (South 27.2, Midwest 13.1, West 8.9). 2014--2018 are in the historical Excel files (about \$1.6bn annual rate in January 2014). Monthly annual rate: \$40bn (April 2025), \$51bn (April 2026), about \$68--75bn (June--July 2026, $+57\%$ y/y; secondary reports of the 1 Sept 2026 release). An industry source's ``\$162bn of BEA data center investment in 2025'' is not this line and should not be used.
  \item \textbf{Equipment in NIPA 5.5.5U}: computers and peripheral equipment (B935RC) \$160.3bn (2022), 149.1 (2023), 178.7 (2024), 270.7 (2025); quarterly annual rate 222.8, 253.8, 280.9, 325.4 in 2025 and 384.2, 400.6 in 2026Q1--Q2 (FRED, 4 Sept 2026 vintage; real computers investment B935RX fell from 356.1 to 348.4 between the two quarters while the nominal series rose). Communication equipment (C275RC) 141.9 (2024), 163.7 (2025). Information processing equipment (Y034RC) 501.1 (2024), 626.6 (2025), 766.5 in 2026Q1. Software 754.8 (2025). Shares of private nonresidential fixed investment (\$4{,}250.8bn, 2025): IPE 14.7\%, software 17.8\%, together 32.5\%. Servers and GPUs are in ``computers and peripheral equipment'' (fixed-asset types EP1A mainframes, EP1B PCs, EP1G storage, EP1H system integrators); networking in EP20; there is no GPU or data center equipment line.
  \item \textbf{Detailed fixed assets by industry $\times$ asset type} (\texttt{detailnonres\_inv1/stk1/dep1.xlsx}, \url{https://apps.bea.gov/national/FA2004/Details/Index.htm}, updated 26 Sept 2025, 1901--2024): \emph{no} data center asset type; structures are SOO1 office, SU20 communication, SC02 other commercial, etc. Industry 5140 in 2024, \$bn: total investment 282.1, of which equipment 78.8 (mainframes 10.4, PCs 11.0, system integrators 19.5, storage 1.0, communications 18.4, other electrical 4.0, general industrial 1.7), structures 36.9 (office 20.7, communication 12.9), IPP 166.3 (software 133.5). Computers (EP1A--H) are 45.2bn, 26\% of the all-industry total (171.4bn); office structures 20.7bn, 21\% of the total (99.6bn). Industry 5110 (Microsoft) invests 12.0bn in computers and 4.4bn in communications equipment; real estate 5310 invests 5.5bn in computers and 15.1bn in office structures.
  \begin{center}\footnotesize
  \begin{tabular}{lrrrrrr}
  \toprule
  Industry 5140 investment, \$bn & 2014 & 2017 & 2019 & 2021 & 2023 & 2024 \\
  \midrule
  Equipment & 22.3 & 37.4 & 37.4 & 59.2 & 70.8 & 78.8 \\
  \quad computers EP1A--H & 14.2 & 19.3 & 20.3 & 34.0 & 37.9 & 45.2 \\
  \quad communications EP20 & 5.7 & 10.8 & 10.0 & 14.5 & 18.5 & 18.4 \\
  \quad other electrical EO72 & 0.6 & 1.8 & 1.8 & 2.6 & 3.8 & 4.0 \\
  Structures & 7.5 & 9.8 & 23.4 & 29.2 & 36.8 & 36.9 \\
  \quad office SOO1 & 3.5 & 4.9 & 13.2 & 16.9 & 20.8 & 20.7 \\
  \quad communication SU20 & 2.8 & 4.2 & 8.4 & 10.2 & 12.7 & 12.9 \\
  Intellectual property (software) & 36.5 & 68.8 & 79.0 & 118.6 & 156.4 & 166.3 \\
  Total & 66.3 & 116.0 & 139.8 & 207.0 & 264.0 & 282.1 \\
  \bottomrule
  \end{tabular}
  \end{center}
  \item \textbf{Net stocks (current cost, 2024)} for 5140: equipment 245bn (communications 84, system integrators 41, PCs 22, mainframes 21, other electrical 14, general industrial 7), structures 304bn (office 171, communication 104), IPP 378bn. Depreciation 2024: equipment 53.2bn (implied rate 22\%), structures 7.4bn (2.4\%), IPP 131.6bn (35\%).
  \item \textbf{Scale check.} Big-4 hyperscaler capex was about \$210--230bn in 2024 and \$360--410bn in 2025 (Section~\ref{sec:boom}), worldwide and including finance leases and non-data-center capital; BEA records 45bn of computers for 5140 and 12bn for 5110 in 2024, and 179bn of computers for the whole economy. Part of the gap is foreign capex; part is that BEA's allocation of equipment across industries is, in the readme's words, ``based on judgmental trends''; part is GPUs bought by CoreWeave/Oracle (industry 5110 or 5415) or leased. The 5140 industry-by-asset detail is usable for \emph{shares} of the bundle, not for its level.
  \item \textbf{Fed measure (Brandsaas, Garcia, Kurtzman, Nichols, Zytek, FEDS 2025-109).} Data center investment $=$ NIPA data center structures $+$ (computers investment $-$ its 2015--2022 linear trend). About \$60bn annual rate in 2024Q4, \$180bn in 2025Q4, \$370bn by 2026Q2 (mean nowcast from Dodge project data). Split 30\% structures, 70\% equipment (from 2024 Census and BEA data); equipment import share 80\%; hence 44\% of data center investment is domestic. Contribution to nominal GDP growth 0.17pp (2024), 0.41pp (2025), 0.75--1.1pp (2026) before the import adjustment. Plans: \$7bn a year in 2019--22, \$300bn in 2024, \$900bn in January--September 2025; abandonment 0.33; plan-to-start 11 months; start-to-completion 19 months for a \$1bn project. \url{https://doi.org/10.17016/FEDS.2025.109}
  \item \textbf{GDP by industry}, value added of 518--519 (current \$bn): 315.6 (2020), 369.1 (2021), 416.8 (2022), 471.2 (2023), 529.1 (2024); real (2017\$) 467 (2022), 539 (2023), 608 (2024). Real output of 518 is deflated with the BLS PPI for 518210, which rises 1.4\% a year (2016--2025; matched-model contracts), whereas hedonic cloud prices fell 7--17\% a year over 2009--2016 (Byrne, Corrado and Sichel 2018) and list prices have been flat since 2017; measured real output and TFP of the industry are understated.
\end{itemize}

\section{Energy use of data centers}\label{sec:energy}

\subsection{Electricity consumption}
\begin{center}\footnotesize
\begin{tabular}{p{3.6cm}lrrp{5.6cm}}
\toprule
Source & Year & TWh & \% of US & Notes \\
\midrule
LBNL 2024 (Shehabi et al.) & 2014 & 58 & 1.5 & \\
 & 2018 & 76 & 1.9 & \\
 & 2023 & 176 & 4.4 & servers $\approx 100$~TWh, of which GPU servers $>40$; storage 16 \\
 & 2028 & 325--580 & 6.7--12.0 & 74--132~GW at 50\% utilisation \\
LBNL 2025 update (Smith et al., June 2026) & 2024 & 192 & 4.7 & 2023--24 revised slightly down \\
 & 2028 & 464 & & reference \\
 & 2030 & 649 (521--843) & 11.8 (9.5--15.3) & AI servers 55\% of DC energy; 148~GW \\
IEA Energy and AI (April 2025) & 2024 & $\approx 187$ & & 45\% of 415~TWh world; $+240$~TWh to 2030 \\
EPRI 2026 & 2024 & 177--192 & & 2030: 380--790~TWh, 9--17\% \\
EIA AEO2026 & 2025 & & & servers 7\% of commercial electricity; 22--33\% by 2050 \\
Grid Strategies 2025 & 2030 & & & utility peak $+166$~GW, of which $\approx 90$~GW data centers \\
\bottomrule
\end{tabular}
\end{center}
Denominator (EIA Electric Power Annual Table 2.2): retail sales 3{,}765~TWh (2014), 3{,}859 (2018), 3{,}874 (2023), 3{,}975 (2024); total end use including direct use 4{,}011 (2023), 4{,}110 (2024). Fleet PUE 1.55 (2018) to 1.45 (2024); AI facilities 1.145; infrastructure (cooling, UPS) is 31--36\% of data center electricity. Reports: \url{https://eta-publications.lbl.gov/sites/default/files/2024-12/lbnl-2024-united-states-data-center-energy-usage-report.pdf}; \url{https://www.osti.gov/servlets/purl/3374245}; \url{https://www.iea.org/reports/energy-and-ai}; \url{https://powering-intelligence.epri.com/executive-summary.html}; \url{https://www.eia.gov/outlooks/aeo/}.

\subsection{Kind of energy: physical vs contractual}
\begin{itemize}
  \item \textbf{Grid electricity} is at least 95\% of physical supply in 2024. No official series separates it.
  \item \textbf{Behind-the-meter gas} (2024--): xAI Colossus, Memphis, about 1.5~GW of turbines operating; Meta Hyperion, Louisiana, three Entergy CCGTs of 2.26~GW (ten plants, $>7$~GW announced); OpenAI/Oracle Stargate Abilene, 2.45~GW planned; Crusoe, Microsoft, Amazon projects. Cleanview (mid-2026): about 2~GW online, 2.8--3.2~GW by end-2026, 82~GW announced since 2025. GE Vernova gas turbine backlog 116~GW (20\% data center customers), Siemens Energy 69~GW, Mitsubishi 35~GW; Global Energy Monitor counts 378~GW of US gas plants in development. National on-site gas is 1--2\% of data center energy today; the announced pipeline would make it a first-order channel from data centers to gas demand by 2028--30. \url{https://cleanview.co/reports/behind-the-meter-data-centers}
  \item \textbf{Diesel backup}: large installed standby capacity (Virginia alone about 27~GW nameplate) run 10--30 minutes a month for testing; negligible energy. \textbf{Fuel cells}: Equinix--Bloom 75~MW operating; negligible.
  \item \textbf{Nuclear PPAs} (contractual, grid-delivered): Microsoft--Constellation Three Mile Island 835~MW from 2028; Meta--Constellation Clinton 1{,}121~MW from 2027; Amazon--Talen Susquehanna up to 1.9~GW; Google--NextEra Duane Arnold 615~MW from 2029; Google--Kairos 500~MW by 2035; Meta--Vistra/TerraPower/Oklo up to 6.6~GW by 2035. \textbf{Renewable PPAs}: the four hyperscalers signed 49\% of the 55.9~GW of global corporate clean-energy PPAs in 2025 (BNEF); Equinix reports 96\% renewable coverage on a market basis. In the IO table the node buys electricity from 2211 regardless of these contracts.
  \item \textbf{Fuel behind the electricity} (EIA 2024 generation): US gas 43\%, nuclear 18\%, coal 15\%, wind 10.5\%, hydro 5.6\%, utility solar 5\%. Virginia: gas 59\%, nuclear 28\%, and sales (138~TWh) exceed in-state generation (103) by a third, met by PJM imports; PJM's 2026/27 cleared capacity is 45\% gas, 22\% coal, 21\% nuclear. Texas: gas 52\%, nuclear 7\%, wind and solar large. Arizona: gas 47\%, nuclear 28\%. Ohio: gas 60\%, nuclear 13\%. Georgia: gas 41\%, nuclear 34\%. Dominion's 2024 IRP: demand $+5.5\%$ a year (doubling by 2039, $+0.5\%$ without data centers), 5.9~GW of new gas in 2030--36 in every scenario. For the model this means the marginal electricity bought by data centers is more gas-intensive than the national average: the data center node interacts with the fossil-generation share of the electricity split in \texttt{data\_notes.tex}.
\end{itemize}

\subsection{Electricity bill and price effects}
At EIA average prices (commercial 12.6--12.8 cents/kWh, industrial 8.0--8.1 cents, 2023--24), 176~TWh cost \$14--22bn in 2023 and 192~TWh \$16--25bn in 2024; large data centers pay between the two rates. PJM capacity prices went from \$28.92/MW-day (2024/25) to \$269.92 (2025/26; \$444 in Dominion) and the \$329.17 cap (2026/27, uncapped \$388.57); the market monitor attributes \$21.3bn (45\%) of the last three auctions' \$47.2bn to data center load, much of it forecast rather than built. CMU/NC State: data centers and crypto raise US generation cost 8\% by 2030 and more than 25\% in northern Virginia. IMF WP 25/081 (CGE): US electricity prices $+8.6\%$ under constrained renewables and transmission. Twenty-four states have large-load tariffs (EEI, August 2026). Harvard ELI (Martin and Peskoe 2025) documents the socialisation of interconnection costs in about fifty rate cases. \url{https://www.pjm.com/-/media/DotCom/markets-ops/rpm/rpm-auction-info/2026-2027/2026-2027-bra-report.pdf}; \url{https://www.imf.org/en/publications/wp/issues/2025/04/21/power-hungry-how-ai-will-drive-energy-demand-566304}

\section{Cost structure of a data center node}\label{sec:cost}

\subsection{Evidence}
\begin{center}\small
\begin{tabular}{p{4.2cm}p{9.8cm}}
\toprule
Source & Electricity share and cost structure \\
\midrule
Census SAS, NAICS 518 & Purchased electricity \$980m in 2017 on \$127bn expenses (0.8\%); purchased fuels \$12m (2012). Series discontinued after 2017 (folded into other expenses; AIES 2023 detail not yet retrieved). 2022: revenue 308bn, expenses 197bn, payroll 90.6bn, depreciation 20.0bn. \\
BEA IO 2017, 518200 & Electricity 1.0\% of gross output; intermediates 41\%, compensation 19\%, gross operating surplus 38\% (Section~\ref{sec:io}). \\
Equinix (10-K, IR) & Revenue \$8.7bn (2024), \$9.2bn (2025); cost of revenues 51\%; utilities 28\% of cash operating expenses and 11--16\% of revenue (about 15\% recently), partly passed through. \\
Digital Realty (10-K) & Utilities expense \$1.0bn (2022), 1.47 (2023), 1.33 (2024), 1.43 (2025) on revenue 4.7--6.1bn: 21--27\% of revenue; tenant reimbursements cover 87--94\% of it. \\
Colocation pricing (CBRE H2 2025) & Wholesale asking rent \$196/kW-month (Ashburn $>\$215$); electricity at 7 cents/kWh and full load is \$51/kW-month, i.e.\ 21\% of the customer's all-in cost or 26\% if bundled into rent. \\
Epoch AI (May 2026), 1~GW AI campus & Capex \$38bn (servers 56--60\%, facility \$11.4/W $\approx 30\%$, network \$4.9/W $\approx 13\%$). Annual TCO \$8.5bn: electricity \$0.6bn (7\%; PUE 1.14, utilisation 71\%, 8.3 cents/kWh), other opex \$0.3bn (4\%), capital services \$7.6bn (89\%), of which servers \$5bn. IT life 5 years, facility 14. \\
Facility only (shell $+$ mechanical/electrical, no IT) & \$10--12m/MW conventional, \$15--20m/MW AI-ready (Goldman, JLL: shell-and-core \$11.3m/MW). Electricity \$0.5--0.6m/MW-year against capital services \$1.4--1.5m and labour/maintenance \$0.2m: electricity 25--30\%, other intermediates $\approx 10$\%, labour small, capital 60--65\%. \\
McKinsey (April 2025), 2025--30 AI stack & \$5.2tn: 60\% chips and hardware, 25\% power, cooling and electrical, 15\% land, site and shell. \\
\bottomrule
\end{tabular}
\end{center}
\url{https://epoch.ai/data-insights/ai-datacenter-cost-breakdown}; \url{https://www.goldmansachs.com/insights/articles/tracking-trillions-the-assumptions-shaping-scale-of-the-ai-build-out}; \url{https://www.mckinsey.com/industries/technology-media-and-telecommunications/our-insights/the-cost-of-compute-a-7-trillion-dollar-race-to-scale-data-centers}

\subsection{Proposed expenditure shares of the node}
The share of electricity depends on what the node's output is. Three consistent definitions:
\begin{enumerate}
  \item \textbf{Computing services, all-in} (output $=$ annual cost of the whole US fleet including IT; the natural definition for a node whose capital includes servers): electricity 7--10\% of gross output, other intermediates (telecom, maintenance, software, water) 5--8\%, labour 3--5\%, capital services 75--85\%. Gross output for 50~GW at Epoch's \$8.5m/MW-year for AI and less for conventional capacity is of the order of \$250--350bn in 2025, consistent with SAS 518 revenue (\$375bn in 2024, which includes non-data-center processing).
  \item \textbf{Facility services} (output $=$ rent of powered shell, IT owned by users, as in colocation): electricity 20--30\%, other intermediates 10\%, labour 5\%, capital 55--65\%.
  \item \textbf{Industry 514 as published}: electricity 1\%. Not a data center.
\end{enumerate}
Recommendation: definition 1 with electricity at 8\% (sensitivity 5--15\%), natural gas 0--2\% (rising with on-site generation), and the capital bundle of Section~\ref{sec:design}. Note that the node's electricity is at most 5\% of US electricity: even a doubling of its demand is a 5\% shift of the electric power row, so the aggregate effects of the boom on electricity run through the level of investment in generation, not through the intermediate share.

\subsection{Depreciation}
Hyperscalers extended server lives from 4 to 6 years in 2022--23 (Microsoft, Alphabet, Amazon; Meta to 5.5 in 2025); Amazon moved part of its fleet back to 5 years in 2025, citing AI; Burry (November 2025) argues 2--3 years. Goldman calls the 4--6 year silicon life the most important scale assumption; Epoch annualises at 5 years (3 years raises the annual cost of a GW by 50\%). Tax: 5-year MACRS for servers, 39-year for the shell, 100\% bonus depreciation after January 2025. BEA (\url{https://apps.bea.gov/national/pdf/BEA_depreciation_rates.pdf}) has no single geometric rate for computers (Oliner used-price profiles); the 0.3119 often quoted is office and accounting equipment; communication equipment 0.11--0.15; office buildings 0.0247 (life 36). Implied rates in the 5140 detailed accounts: equipment 22\%, structures 2.4\% (2024). For the node: IT $\delta$ 0.25--0.40 (3--6 years), mechanical/electrical 0.08--0.15 (10--20 years), shell 0.025 (36 years); bundle-weighted at the 60/25/15 capex mix this gives $\delta \approx 0.20$--0.28, far above any existing sector.

\section{Construction of the node in the 41-sector dataset}\label{sec:design}

\subsection{Which sector to split}
\begin{itemize}
  \item \textbf{Option A: split \texttt{Info} (26) into Data centers and Rest of Info.} Consistent with the classification: hosting, cloud infrastructure and colocation are 518210; hyperscaler own-account capacity is capital of 514/511, both inside \texttt{Info}. Data on the 518200 column and row exist for 2007/2012/2017; investment by asset exists annually for 5140. Rest of Info keeps publishing, software, telecom, broadcasting, and internet publishing.
  \item \textbf{Option B: split \texttt{Real\_estate} (28).} Not supported: BEA records no colocation output in real estate; Equinix and Digital Realty are REITs by tax status only; the real estate sector is dominated by housing services. What real estate does contain is landlords' purchases of electricity for tenants (\$41bn in 2023, of which data center landlords are a part) and 15bn a year of office structures investment, some of it data centers built by developers for lease. Treatment: reassign an estimated data center share of ORE's electricity purchases and of its office structures investment to the node.
  \item \textbf{Functional definition.} Follow the PwC/Data Center Coalition and Fed practice and define the node as all US data center capacity (hyperscale, colocation, enterprise) regardless of owner, selling computing services to the sectors that use it, including its former parents. In the model this reassigns capital from \texttt{Info}, \texttt{Finance}, \texttt{Retail}, etc.\ to the node and adds an intermediate purchase in return; it is an accounting reallocation with no effect on the aggregate.
\end{itemize}

\subsection{Construction by object}
\begin{center}\footnotesize
\begin{tabular}{p{2.2cm}p{3.5cm}p{3.6cm}p{3.0cm}p{2.0cm}}
\toprule
Object & Source & Construction & Assumption & Status \\
\midrule
Gross output, snapshot 2023/24 & FERC 50~GW; LBNL 192~TWh; Epoch/JLL cost per MW; SAS 518 revenue & capacity $\times$ annual cost per MW by type (AI vs conventional); cross-check with SAS hosting revenue plus hyperscaler own-use & steady state is a 2023--24 snapshot; no time series & feasible \\
IO column & 2017 detail 518200 column; Epoch/facility cost shares; EIA prices $\times$ TWh for electricity & electricity from 2211 at 8\% (5--15\%); gas from 2212/211 at 0--2\%; telecom from Rest of Info; maintenance and software from Prof\_Tech and Rest of Info; labour and GOS residual & 518200 non-energy structure re-scaled to the engineering shares & feasible \\
IO row & 2017 detail 518200 row; SAS/Census product lines; IDC cloud revenue by customer industry & sales to Rest of Info (hyperscaler own use), Prof\_Tech, Finance, Retail, Mgmt, government, PCE (consumer cloud) & drop payment-processing customers; proportional to IT spending by industry & moderate \\
Investment column ($\gamma^I$) & FA 5140 by asset 2019--2024; NIPA 5.4.5U data centers; Census C30 & bridge: SOO1, SU20 $\to$ Construction; EP1A--H, EP20 $\to$ Computers (334); EO72 (UPS, switchgear) $\to$ Electrical (335); EI50 (cooling) $\to$ Machinery (333); software $\to$ Info/Prof\_Tech; shares 60/25/15 IT/M\&E/shell & vom Lehn--Winberry bridge on 5140 shares; level from NIPA and Fed measure & feasible \\
Capital stock, $\delta$ & FA 5140 net stocks; asset lives above & bundle-weighted $\delta \approx 0.2$--0.28 & 5140 stock scaled to functional capacity & feasible \\
Import content & Fed: 80--90\% of IT equipment imported & see below & model choice & open \\
Shock & none of its own & recipient of the boom shock (Section~\ref{sec:boom}) & & open \\
Other sectors & remove reassigned capital and add intermediate purchases from the node & keep 1997--2023 series for the 40 other sectors & & feasible \\
\bottomrule
\end{tabular}
\end{center}

\subsection{The import problem}
In RbcProdNet the investment good of every sector is produced domestically; sector 9 \texttt{Computers} would produce the servers. In the data 80--90\% of the high-tech equipment for the buildout is imported (Fed FEDS 2025-109 and Fiori, Lipa and Nuenninghoff 2026), and the Fed's current-account note treats the boom as an investment-specific technology shock that widens the deficit by 0.2\% of GDP. Options: (i) keep the closed economy and read the Computers sector's response as world supply; (ii) add an import leak for the IT part of the data center bundle at an exogenous world price, so that the domestic content of the boom is structures, electrical and cooling equipment, and construction; (iii) treat the IT bundle as a non-produced input in fixed supply (a resource constraint, like the fuel rows in \texttt{data\_notes.tex}). Option (ii) is the honest one and changes the aggregate multiplier of the boom by a factor of about two. The plan adopts (ii) with the import row at an exogenous world price and balanced trade, combined with the US+Nvidia addition to sector 9 (Section~\ref{sec:nowcast}), so that the domestic content of the IT bundle is 60\% gross rather than the Fed's 20\%.

\section{Modelling the investment boom}\label{sec:boom}

\subsection{Facts}\label{sec:facts}
\begin{itemize}
  \item \textbf{Hyperscaler capex} (cash purchases of property and equipment, \$bn; definitions differ across firms, leases excluded unless noted): Amazon 63.6 (2022), 52.7 (2023), 83.0 (2024), 131.8 (2025), about 220 guided for 2026; Alphabet 31.5, 32.3, 52.5, 91.4, 195--205; Meta 31.2, 27.0, 37.3, 69.7, 130--145; Microsoft (fiscal years) 23.9, 28.1, 44.5, 64.6, about 175 calendar 2026 including leases; Oracle 4.5, 8.7, 6.9, 21.2, 55.7 (FY26); CoreWeave 35--39 (2026). Big-4 total about 210--230 (2024), 360--410 (2025), 720--750 guided for 2026. Alphabet: 60\% of 2025 capex servers, 40\% data centers and networking.
  \item \textbf{Industry totals}: Dell'Oro global data center capex \$455bn (2024, $+51\%$), $+57\%$ in 2025 (about \$715bn), above \$1tn in 2026; McKinsey \$6.7tn cumulative to 2030 (\$5.2tn AI); Goldman \$7.6tn over 2026--31 (\$765bn in 2026, \$1.6tn in 2031; assumes \$15m/MW facilities, \$2{,}500/kW new power, 4--6 year silicon life); BIS: AI investment about 1\% of GDP in the most exposed economies.
  \item \textbf{GDP accounting}: information processing equipment plus software contributed 1.5pp to annualised real GDP growth in 2025H1 with the rest of GDP at $-0.1$pp (Furman's ``92\% of growth''); IPE alone 0.90pp in 2025Q1 and 0.16pp in 2025Q3 (St.\ Louis Fed); net of imported computers, semiconductors and telecom equipment the AI contribution falls to 0.4--0.5pp (MRB) and Goldman puts the data center boom at 0.2pp of 2025 growth; the Fed's project-based measure gives 0.4pp in 2025 and 0.75--1.1pp in 2026 before imports; Apollo/Slok: data center capex 0.6\% of GDP (2023), 1.4\% (2025), 3.1\% (2027); Bridgewater 1.4pp of growth in 2026. IP and equipment investment as a share of GDP is 0.8pp above 2024 and near its 2000 peak of 11.5\% (Li 2026).
\end{itemize}

\subsection{How the literature introduces the boom}
\begin{center}\footnotesize
\begin{tabular}{p{3.6cm}p{4.2cm}p{6.4cm}}
\toprule
Paper & Shock & Result \\
\midrule
Fiori, Lipa, Nuenninghoff, FEDS Note 14 July 2026 & Investment-specific technology shock (Fisher 2006 SVAR, 1993--2019); the AI boom as its current phase & one s.d.: investment $+2\%$, output $+0.5$pp, current account $-0.2\%$ of GDP; 90\% of high-tech equipment imported from East Asia \\
Li, FEDS Note 6 July 2026 (Li 2007) & Embodied general-purpose technology of unknown size; learning by investing; rational overshooting & railroads, canals, 1990s fiber as analogues; investment can overshoot without irrationality; do not curb investment pre-emptively \\
IMF WP 25/081 \emph{Power Hungry} (2025) & Exogenous path of data center electricity demand in a global CGE & US electricity prices $+8.6\%$ with constrained renewables and transmission; US CO$_2$ $+5.5\%$ \\
Epoch GATE (Erdil et al.\ 2025) & Compute stock automates tasks; integrated assessment & compute build-out first, TFP later \\
Kleinman Center (Penn) 2026 & Nodal power-flow GE & a 200~MW data center at a constrained node can double the local electricity price \\
Acemoglu, \emph{Economic Policy} 2025 & Task-based TFP benchmark & TFP $+0.66\%$ over ten years (conservative size for a news shock) \\
Korinek and Suh, NBER w32255 (2024) & Automation of a task continuum & GDP $+100$--300\% over ten years (upper tail) \\
Beaudry and Portier 2006; Jaimovich and Rebelo 2009 & News about future TFP; investment-specific news & the standard way to get an investment boom before measured productivity in a neoclassical model with adjustment costs \\
\bottomrule
\end{tabular}
\end{center}
No paper introduces the boom as a shock to a sector node in a vom Lehn--Winberry investment network. The timing of 2023--26 (investment first, productivity not yet measured) is that of a \emph{news shock}: an anticipated increase in the productivity of sectors that use computing services, or in the productivity of the data center node itself, announced at $t$ and realised (or not) at $t+k$. This is also the setting in which the planner's precautionary motives from RbcProdNet bite: the news may not materialise, the capital is short-lived and specific, and the electricity it needs is supplied by a sector with long-lived capital. Alternatives: an investment-specific technology shock to the IT bundle (fits 2010--2022 falling GPU prices, not 2024--26 chip and memory inflation); a demand shift toward computing services (rhetorical unless justified by productivity); an exogenous investment path in the node matching the Fed's \$370bn by 2026Q2 (useful as accounting and as a robustness check, silent on expectations and on a bust). Existing IO studies of data centers (PwC for the Data Center Coalition, IMPLAN, 2026: 5.5m jobs and \$927bn of GDP in 2024; Virginia JLARC 2024: 74{,}000 jobs, \$9.1bn) treat imported IT as domestic value added and are not usable for calibration; JLARC's finding that most local benefits are in construction and that unconstrained Virginia load doubles in ten years is the relevant part.

\subsection{The shock: compute-intensity shift versus compute-augmenting technical change}\label{sec:shockdesign}

The plan (D9, S4) introduces the boom as a rise in the weight of computing services in every buyer's technology, $\mu_{c,jt}$, rather than as a productivity gain attached to compute. This subsection records why, with the algebra and the references, because the choice is the one most likely to be questioned.

\paragraph{Setup.} Buyer $j$ combines a bundle $B$ of its other inputs (the existing value-added--materials nest) with computing services $C$ in a CES nest,
\[
Q = \Big[(1-\mu)^{1/\sigma} B^{\frac{\sigma-1}{\sigma}} + \mu^{1/\sigma} (A_c C)^{\frac{\sigma-1}{\sigma}}\Big]^{\frac{\sigma}{\sigma-1}},
\qquad
P = \Big[(1-\mu) P_B^{1-\sigma} + \mu\, \tilde p^{\,1-\sigma}\Big]^{\frac{1}{1-\sigma}},
\quad \tilde p \equiv p_c/A_c ,
\]
where $p_c$ is the price of compute, $\tilde p$ its price per efficiency unit, $P$ the buyer's unit cost, and $\sigma$ the elasticity of substitution between compute and everything else. Three candidate shocks: (a) a rise in $A_c$ (compute-augmenting technical change in the buyer: the same compute does more); (b) a fall in $p_c$ (TFP in the Compute sector, or cheaper GPUs passed through), which enters only through $\tilde p$ and is therefore the same shock as (a) from the buyer's side; (c) a rise in $\mu$ (the intensity shift). Cost minimisation gives
\[
A_c C = \mu \left(\frac{\tilde p}{P}\right)^{-\sigma} Q ,
\qquad
s \equiv \frac{p_c C}{P Q} = \mu \left(\frac{\tilde p}{P}\right)^{1-\sigma},
\]
with $s$ the nominal share of compute in the buyer's gross output, the object we observe (S4.2).

\paragraph{(a)/(b): what a cheaper or better compute input does.} Differentiating with $P_B$ and $p_c$ fixed, using $d\ln P = s\, d\ln\tilde p$ and writing $\varepsilon_j$ for the demand elasticity of $j$'s own output (so $d\ln Q = -\varepsilon_j\, d\ln P$),
\[
\frac{d\ln s}{d\ln A_c} = -(1-\sigma)(1-s),
\qquad
\frac{d\ln C}{d\ln A_c} = -(1-\sigma) + s\,(\varepsilon_j - \sigma).
\]
Two conclusions. First, the nominal share of compute falls whenever $\sigma<1$ and rises only if $\sigma>1$, independently of the output response: this is Hicks's (1932) result that factor-augmenting technical change is biased \emph{toward} the factor if and only if the elasticity of substitution exceeds one, restated in Acemoglu (2002, \emph{REStud}, ``Directed technical change'', and 2003, \emph{JEEA}, ``Labor- and capital-augmenting technical change''), and it is the mechanism behind Baumol's cost disease in Ngai and Pissarides (2007, \emph{AER}) and Herrendorf, Rogerson and Valentinyi (2014, \emph{Handbook}): the sector whose price falls fastest loses nominal share when goods are complements. Second, the \emph{physical} quantity of compute falls when $\sigma<1$ and the compute share $s$ is small: fewer efficiency-adjusted units are needed to do the same job, and the output expansion $s(\varepsilon_j-\sigma)$ is second order because $s$ is a few percent of gross output. So with complements a compute-augmenting shock predicts \emph{less} compute bought, a \emph{smaller} dollar share and a \emph{lower} price of compute. The data show the opposite on all three counts for 2023--25 (Section~\ref{sec:facts}: data-center investment from 0.6\% to 1.4\% of GDP, Nvidia's margin, HBM and packaging prices up). Measuring compute in efficiency units does not rescue the story: $s = p_c C/(PQ)$ is invariant to units, and the sign of $d\ln s$ is $-(1-\sigma)$ regardless of how $C$ is measured. Aghion, Jones and Jones (2019, in \emph{The Economics of Artificial Intelligence}, NBER) make the same point for AI: with $\sigma<1$ the automated tasks' share of GDP shrinks as they get cheap, and growth is bottlenecked by what is not automated.

\paragraph{Is $\sigma<1$ the right prior?} For intermediates across sectors Atalay (2017, \emph{AEJ: Macro}) estimates elasticities near zero (0.0--0.2), and the base model's own $\sigma_q$ is below one; for capital--labour at the plant level Oberfield and Raval (2021, \emph{Econometrica}) find 0.5--0.7. The cases for $\sigma>1$ are IT-specific: Krusell, Ohanian, R\'{\i}os-Rull and Violante (2000, \emph{Econometrica}) estimate 1.67 between equipment and unskilled labour (0.67 with skilled), Eden and Gaggl (2018, \emph{RED}) find ICT capital substitutes for routine labour, and Karabarbounis and Neiman (2014, \emph{QJE}) need $\sigma\approx1.25$ for cheaper investment goods to explain the falling labour share. The 1990s IT boom is the relevant precedent (Jorgenson 2001, \emph{AER}; Tevlin and Whelan 2003, \emph{JMCB}): computer prices fell 15--20\% a year while the nominal IT investment share of GDP roughly doubled, which requires $\sigma>1$ if the shock was price alone. That episode is why $\sigma_c=2$ is a genuine sensitivity in the plan and not a formality. But the identification for 2023--25 is different from the 1990s: then the price fell and the share rose, which is compatible with either an elastic demand curve or a shift; now the price \emph{rose} and the share rose, which no movement along a demand curve produces at any $\sigma$. A rise in both price and quantity is a demand shift by revealed preference.

\paragraph{(c): the intensity shift.} Differentiating in $\mu$ at fixed prices, $d\ln(A_c C) = d\ln\mu + \ldots$ and $d\ln s = d\ln\mu + O(s)$; at $\sigma=1$ exactly $s=\mu$. With an upward-sloping short-run supply of compute (adjustment costs on GPU\_servers, DC\_facilities and Compute) the price of compute rises as well. Quantity up, share up, price up, for any $\sigma$: it is the only one of the three shocks that matches the three facts jointly, and at $\sigma=1$ it is identified directly from the expenditure shares without a compute deflator, whose candidates differ by an order of magnitude (BEA computers deflator versus Epoch FLOP/\$; plan S4.4). A supply-side reading of the same object: the planner's shadow price of compute rises because compute has become more useful at the margin at the same efficiency units, which is what ``more tasks can be done with compute'' means.

\paragraph{Why this is Acemoglu's framing.} In the task model of Acemoglu and Restrepo (2018, \emph{AER}, ``The race between man and machine''; also Acemoglu and Autor 2011, \emph{Handbook}) output is a CES aggregate over a continuum of tasks, tasks below a threshold $I$ are performed by capital and the rest by labour, and the aggregate collapses to a two-factor CES whose \emph{weights} are the measures of the two task sets:
\[
Y = \Big[(I-N+1)^{1/\sigma}(A_K K)^{\frac{\sigma-1}{\sigma}} + (N-I)^{1/\sigma}(A_L L)^{\frac{\sigma-1}{\sigma}}\Big]^{\frac{\sigma}{\sigma-1}} .
\]
Automation is $dI>0$, a change in the weights, and they show it always raises capital's share and displaces labour at given factor prices, whereas capital-augmenting change $dA_K$ raises labour's share whenever $\sigma<1$ and has, in their words, counterintuitive implications for the factor it augments; Acemoglu and Restrepo (2018, \emph{AEA Papers and Proceedings}, ``Modeling automation'') is the short statement of exactly this contrast and the reason automation cannot be written as factor-augmenting technical change. The lineage is Zeira (1998, \emph{QJE}) for tasks moving from labour to capital and Acemoglu and Autor (2011) for the assignment; the AI application is Acemoglu and Restrepo (2019, in \emph{The Economics of Artificial Intelligence}, NBER). Their productivity effect is a separate term. Acemoglu (2025, \emph{Economic Policy}, ``The simple macroeconomics of AI'') carries the same structure to AI: the aggregate effect is the fraction of tasks exposed (a shift in the task set, i.e.\ in the weights) times the cost saving per task, with the TFP consequence by Hulten's theorem, 0.66\% over ten years. Our $\mu_{c,jt}$ is the task-set shift for compute in each buyer; the add-on $A^c_t$ (plan S4.7) is the cost saving per task, scaled to his 0.66\%. Keeping them separate is the Acemoglu--Restrepo decomposition into displacement (here: reallocation toward compute) and productivity, and it lets the paper say how much of the boom is the first without asserting anything about the second. The demand-side analogue in the structural-change literature is Comin, Lashkari and Mestieri (2021, \emph{Econometrica}), where nonhomothetic CES weights shift with income, and Hubmer (2023, \emph{REStud}), who separates preference shifts from relative-price effects; our shifter plays the same role on the technology side.

\paragraph{Units, and the effect of the shift at fixed inputs.} With the CES written as above, a change in $\mu$ at fixed inputs is not output-neutral: at $\sigma=1$, $d\ln Q=\ln(C/B)\,d\mu$, whose sign depends on the units in which $C$ and $B$ are measured. One option is to remove the term by the normalised CES of Klump and de La Grandville (2000, \emph{AER}) and Klump, McAdam and Willman (2012, \emph{JES}), with the 2018 point as normalisation. The task model says this would be a mistake: the term is the productivity effect, and the right response is to pin the units, not to normalise the term away (next paragraphs).

\paragraph{Reduced form of a task model: does the shift imply a productivity gain?} The CES with task-set weights is an exact aggregation of the task continuum at every date, not a limit (Acemoglu and Restrepo 2018, Proposition 1, eq.~12, pp.~1498--99: ``the share parameters are endogenous and depend on the state of the two types of technologies and the equilibrium choices of firms''; Corollary 1, p.~1500, is the Cobb--Douglas case with the automated task share as exponent), but its weights are not parameters: they are the outcome of an assignment. What is exogenous is the frontier $I$ (which tasks \emph{can} be done by capital); adoption $I^*=\min\{I,\tilde I\}$ is endogenous, and a task is reassigned only if capital does it more cheaply. The productivity gain is not assumed; it is their Proposition 3 (pp.~1501--03), where productivity is defined as output at fixed $K$ and $L$ and
\[
d\ln Y|_{K,L} = \frac{B^{\hat\sigma-1}}{1-\hat\sigma}\Big[\big(W/\gamma(I^*)\big)^{1-\hat\sigma} - R^{1-\hat\sigma}\Big]\,dI ,
\]
``the cost savings from automation, which are given by the difference between the effective wage at $I^*$ and the rental rate'' (p.~1503); in the unconstrained case ($I^*=\tilde I<I$) the marginal task is indifferent and ``additional automation technologies do not'' raise productivity (p.~1502). So the gain has the sign of the shift but no fixed size: it is the cost gap at the reassigned tasks, zero at the margin (``so-so automation'', Acemoglu and Restrepo 2019, \emph{JEP}) and large when the frontier jumps discretely, as with tasks that were infeasible and become feasible well below labour cost. It is measured as Hicks-neutral TFP of the final good, but there is no separate productivity factor: it is the derivative of the aggregate function with respect to its weights, and the same weight shift moves marginal products asymmetrically (the displacement term). Acemoglu (2025) reduces this to one formula, TFP gain $=$ share of tasks exposed $\times$ average cost saving per task, Hulten applied to the reassigned tasks. The cost saving itself is not estimated in the 2018 paper; Acemoglu and Restrepo (2022, \emph{Econometrica}) impose one for industrial robots (about 30\%, from the robot-adoption evidence; check the figure before citing), and Acemoglu (2025) measures it for generative AI from task-level experiments (Noy and Zhang 2023; Brynjolfsson, Li and Raymond 2023; Peng et al.\ 2023): about 27\% on the tasks studied, about 14\% after his adjustment for hard tasks, gross of the cost of the compute that does the task.

\paragraph{The cost saving in our nest: task units.} In A--R the unit cost index is $P^{1-\sigma}=\int p(i)^{1-\sigma}di$ and moving a measure $dI$ of tasks from labour to capital changes it by $[R^{1-\sigma}-(W/\gamma(I))^{1-\sigma}]\,dI$; the saving is the per-task price gap between the two ways of doing the same task. Our unit cost is $P_j^{1-\sigma}=(1-\mu)P_B^{1-\sigma}+\mu\,\tilde p_c^{\,1-\sigma}$, and
\[
d\ln P_j=-\,\frac{P_B^{1-\sigma}-\tilde p_c^{\,1-\sigma}}{(1-\sigma)\,P_j^{1-\sigma}}\,d\mu\;\xrightarrow{\ \sigma\to1\ }\;-\ln\!\frac{P_B}{\tilde p_c}\,d\mu .
\]
The unnormalised CES already generates the productivity effect, with the sign of $P_B-\tilde p_c$: the shift lowers cost iff a task done with compute is cheaper than a task done with the bundle. What is undetermined in a dollar CES is $\tilde p_c/P_B$ itself, because at $\sigma=1$ the share $s=\mu$ carries no information on per-task prices. We therefore fix the units of $C$ by a calibration: write $C_j/\bar c_j$ in the nest and choose $\bar c_j$ so that at 2018 prices one task done with compute costs $(1-\pi)$ times one task done with the bundle, $\tilde p_{c,j,2018}/P_{B,j,2018}=1-\pi$. Then at $\sigma=1$ the shift lowers the buyer's unit cost by $\ln(1-\pi)\,d\mu\approx-\pi\,d\mu$: Acemoglu's formula, tasks moved times saving per task, with $\pi$ a structural parameter (the per-task relative price at the base date) identified from the task-level experiments, $\pi\in\{0,\,0.14,\,0.27\}$, and $\pi=0$ so-so automation. The moving exogenous object is the frontier $\mu_{c,jt}$, as in A--R; adoption at the frontier is rational because $\pi>0$; and the technology-constrained regime we are implicitly in requires $\tilde p_{c,jt}<P_{B,jt}$ at every date, which the simulation checks rather than assumes. Relative to a fixed Hicks-neutral productivity term tied to $\Delta\mu$, this makes the saving a relative price with its own dynamics: the 2023--25 rise in $p_c$ shrinks the saving on newly reassigned tasks, and $q_t$ (cheaper compute per efficiency unit) enlarges it later in the decade, which is A--R's statement that the gains depend on the cost gap at current prices. It also disposes of the normalised CES: the normalisation was a device to kill a term we now calibrate. The cost is that the productivity effect is no longer a separate object but sits inside the buyer's cost index, jointly with $\mu$, $\pi$, $p_c$ and $q_t$; the paper's decomposition is therefore the model at $\pi=0$ against $\pi>0$. Identification of $\mu$ from shares is unchanged at $\sigma_c=1$; at $\sigma_c\neq1$ the per-task relative price enters the share formula through $\pi$, and $\mu$ is re-derived for each $(\sigma_c,\pi)$.

\paragraph{Microfoundation, not model.} The nest is adopted as the technology-constrained regime of the A--R task aggregate: buyers combine tasks with elasticity $\sigma_c$, a measure $\mu_{c,jt}$ of them can be done with compute and all of those are, because compute is cheaper than the bundle at the marginal feasible task. In that regime the aggregate is exactly the CES above with the frontier as weight and the cost gap as productivity effect, so the task framework adds language and two definitions ($\mu$ is the compute-feasible task share, $\pi$ is the cost gap) and no equation. The full task model would add only the other regime, endogenous adoption short of the frontier when compute is too expensive at the margin (A--R's $\sigma_{\text{free}}>\hat\sigma$). It is a real mechanism, but it needs a comparative-advantage schedule $\gamma_j(i)$ per buyer that is unobservable, pins units of compute per task that no data give, and introduces the kink $I^*=\min\{I,\tilde I\}$ into a 40-sector perfect-foresight transition with expectation revisions; and it changes nothing on the network, energy or import side. One choice the framework forces and the reduced form leaves vague: what compute displaces. In A--R it is labour, and the results are about the labour share. Our nest puts compute against the whole bundle $B$ (value added and materials), so the displacement falls on labour, capital and purchased inputs in proportion. For the network questions that is the right choice (AI compute substitutes for call centres, software, professional services and in-house IT as much as for payroll, and the propagation of that displacement through the materials rows is part of the paper), but readers who see the A--R citation will expect labour, so it is stated as a choice; moving the nest inside value added against $L$ is a small code change and a different paper.

\paragraph{What the choice costs.} With $\pi=0$ the intensity shift has no welfare content of its own: it says firms want more compute, not that compute makes them more productive, and the paper cannot claim a GDP gain from the shift itself, only from the investment response and the reallocation. With $\pi>0$ the gain is generated by the reassignment at current prices and disappears with it in the reversal, which is what a task model implies. The comparison shock in the plan (Section 5, item 2(v): compute-augmenting TFP news of Acemoglu's size, with no shift) is there to show that the alternative, with $\sigma<1$, produces a shrinking compute share and cannot generate the observed boom at all.

\section{Nowcast inputs for 2025 and the US+Nvidia treatment}\label{sec:nowcast}
The execution plan is \texttt{data\_center\_nowcast\_plan.md} in this folder. This section records the facts it relies on, verified on 4 September 2026.

\subsection{Vintage of the official network}
\begin{itemize}
  \item Annual IO (71 industries, after redefinitions), GDP by industry and fixed assets run through \textbf{2024} (releases of 25--26 September 2025). The 2026 annual update, the first with national, industry and regional accounts on the same day, is \textbf{30 September 2026} (revision window 2021Q1--2026Q1; first use of Census AIES 2023; no change to computer service lives; no data-center asset type; no 2022 detailed benchmark announced). Until then 2025 industry accounts are a nowcast. \url{https://www.bea.gov/information-updates-national-regional-economic-accounts}
  \item The static zip of all IO tables is no longer served reliably (HTTP~500 on 4 Sept 2026); the 1997--2023 files are saved locally, and 2024 comes from the API (dataset \texttt{InputOutput}, integer \texttt{TableID}, free key) or from the iTables export. The 138-industry GVC Make/Use/import matrices 2007--2024 (May 2026) are the only official annual tables finer than 71 industries.
  \item Quarterly GDP by industry: 2026Q1 (25 June 2026). NIPA underlying detail tables 5.4.5U, 5.5.5U and 5.6.5U: 2026Q2 (26 August 2026 second estimate). Census C30: July 2026 (1 Sept 2026). QSS 518: 2026Q2. The vom Lehn--Winberry replication network (Harvard Dataverse, \texttt{doi:10.7910/} \texttt{DVN/CALDHX}) stops in 2018; no public extension exists; BEA's last capital-flow table is 1997.
\end{itemize}

\subsection{Nowcast of data-center investment, 2025 and 2026H1}
\begin{center}\footnotesize
\begin{tabular}{p{6.6cm}rrrp{3.4cm}}
\toprule
\$bn, nominal & 2023 & 2024 & 2025 & 2026H1, annual rate \\
\midrule
Structures: Census C30 private data centers (annual, latest vintage) & 20.0 & 34.8 & 49.7 & 65--75 (May--July SAAR 65.7, 70.8, 75.2) \\
Computers B935RC (NIPA 5.5.5) & 149.1 & 178.7 & 270.7 & 392.4 (Q1 384.2, Q2 400.6) \\
\quad 2015--2022 linear trend ($92.8 + 8.47\,t$) & 160.6 & 169.1 & 177.6 & 186.0 \\
\quad residual $=$ Fed ``data-center equipment'' & $-11.6$ & 9.6 & 93.2 & 206.4 \\
Communication equipment C275RC & 142.9 & 141.9 & 163.7 & \\
Fed NIPA method: structures $+$ residual & 8 & 44 & 143 & 270--280 \\
Fed Dodge nowcast, mean (FEDS 2025-109, Table 2, Q4 annual rate) & 10 & 58 & 179 & 370 (2026Q2); 464 (2026Q4) \\
\bottomrule
\end{tabular}
\end{center}
The Fed's Table 2 values are fourth-quarter annual rates, not annual totals: a path running from 58 (2024Q4) to 179 (2025Q4) averages about 140 over 2025, consistent with the NIPA-method 143. Annualised, the Fed path is roughly 10 / 40 / 140 / 400bn for 2023--26. The Fed measure is an increment (structures at level, computers in excess of trend) and excludes owner-furnished electrical and mechanical equipment, which is in neither C30 nor the computers residual. 2027 scenarios (Table 2, mean, Q4 annual rate): 357 / 612 / 921bn for plan inflows at 0.25$\times$ / 1$\times$ / 2$\times$ the 2024--25 pace; import adjustment 44\% domestic (30\% structures, 70\% equipment at 80\% imported). The 2023 residual is negative on the current vintage: the detrending is only usable from 2024. Communication equipment was flat through 2024 and rose \$22bn in 2025, an upper bound for the networking part of the bundle. Hyperscaler cash capex 2025: Amazon 128.3, Alphabet 91.4 (60\% servers, 40\% data centers and networking), Meta 69.7 (72.2 with finance leases), Microsoft FY25 64.6 (about 88 with finance leases); 2026 guidance 720--745 for the four; Oracle FY26 55.7; CoreWeave 2026 capex 35--39. QSS NAICS 518 revenue 2025: \$422bn (upper bound for the node's output; it includes payment processing and non-data-center hosting). Capacity end-2025: about 50~GW online (FERC), 35~GW under construction (JLL). QCEW 518210, 2024: 484{,}000 jobs at \$194{,}800 average pay, \$94bn of wages; CES 518 employment fell from 481k to 473k during 2025 while capex tripled, so the node's labour must not be scaled with its capital.

\subsection{Where Nvidia is in the national accounts}
\begin{itemize}
  \item \textbf{Nvidia's accounts} (fiscal year 2026 ended 25 January 2026 $\approx$ calendar 2025; 10-K segment note verified at \url{https://www.sec.gov/Archives/edgar/data/1045810/000104581026000021/R25.htm}): revenue \$215.9bn (FY2025 130.5, FY2024 60.9); Data Center \$193.7bn (compute 162.4, networking 31.4; FY2025 115.2, FY2024 47.5); gaming 16.0. Cost of revenue \$62.5bn (gross margin 71.1\%), R\&D \$18.5bn, SG\&A \$4.6bn, operating income \$130.4bn (60.4\% of revenue), net income \$120.1bn, capex \$6.1bn, 42{,}000 employees. Revenue by customer headquarters: US \$149.6bn (69\%), Taiwan 42.3, China incl.\ Hong Kong 19.7, other 4.3; Nvidia estimates that 76\% of Data Center revenue from Taiwan-headquartered customers (the ODMs) went to end customers in the US and Europe. The older billing-location series (FY2025: US 61.3, Singapore 23.7, Taiwan 20.6, China 17.1) is an invoicing geography (shipments to Singapore under 2\%) and must not be mixed with the HQ series. Long-lived assets: US \$5.1bn, Taiwan 3.2, Israel 1.5. Suppliers: TSMC and Samsung (wafers), SK Hynix, Micron, Samsung (HBM), TSMC CoWoS packaging, Foxconn, Wistron, Fabrinet (assembly). GPU manufacturing cost about \$6.4k for a B200 (HBM 45\%, packaging 17\%, logic die under 15\%) against a \$30--40k price: most of the price is Nvidia's margin. Two direct customers were 22\% and 14\% of FY2026 revenue.
  \item \textbf{Not in US industry 334.} Industry 334 gross output in the summary Use table: \$302bn (2017), 318 (2021), 340 (2022), 342 (2023), while Nvidia's revenue went 27 $\to$ 61bn over 2022--23 and is 216bn in 2025. Nvidia's sales are not in the output of computer and electronic products; a fabless firm's establishments are classified by their activity (design: 5415/5417, or wholesale 4234, or headquarters 55), and the ECPC's 2010 recommendation to classify factoryless goods producers in manufacturing was never implemented (2014 Federal Register notice; no FGP industry in NAICS 2027 either).
  \item \textbf{Epoch AI, 24 August 2026}, \emph{The Nvidia-sized hole in US GDP statistics} (Juniewicz, Carey, Trammell, Ho; \url{https://epoch.ai/publications/the-nvidia-sized-hole-in-us-gdp-statistics}): US GDP growth understated by about 0.3pp over the past year because the import price of an AI server includes Nvidia's markup, which is never recorded as a US export of goods, IP, services or merchanting (the ODM transforms the chips, so it is not merchanting; Nvidia pays TSMC rather than licensing to it); BEA confirmed to them that change-of-ownership data are not collected. Gap about \$30bn a quarter in 2025 (about \$140bn a year for semiconductors); could approach 2pp of growth a year by 2028; GDI already contains Nvidia's US wages and profits (about \$120bn pre-tax in the year to January 2026), which is one reason GDI and GDP diverge. BPM-7/SNA 2025 would record the sales of a factoryless producer as goods exports; BEA has stated since 2014 that it will not do so in the core accounts.
  \item \textbf{The dispute.} The X exchange between Brian Albrecht and the Epoch authors could not be retrieved. The substantive objections on record (Magica, J.P.\ Morgan AM commentary): FGP treatment requires that Nvidia (a US resident) be the economic owner of the wafers and chips and bear the production risk, which needs contract-level evidence; a Taiwanese branch with substantial operations would be a Taiwan resident; the 0.3pp has no firm-level reconciliation; BEA's 2014 ``no GDP impact'' statement concerned a reclassification of recorded flows, not the unrecorded export Epoch describes, so the two are not in contradiction. Related but different measurement points: Korinek and McKelvey (PIIE, May 2026) argue AI output is understated through deflators (inference prices $-94\%$ a year), not through trade; Furman, Hatzius, the St.\ Louis Fed and FEDS Note of 17 July 2026 discuss the import offset to AI investment, not the missing export. For our purposes the dispute matters in one place: Nvidia's operating income is already in NIPA domestic corporate profits (GDI) while the same dollars sit in imports on the expenditure side, so a model economy with ``US+Nvidia'' has a GDP above official GDP by Nvidia's US-installed value added (about \$65--70bn, 0.25\% of 2025 GDP, close to Epoch's 0.3pp of growth); we report that and reconcile it against GDI rather than presenting model GDP as official GDP.
  \item \textbf{Trade side.} US imports of computers (Census end-use 21300) \$88.5bn (2023), 116.8 (2024), 218.2 (2025); computers, peripherals and parts (NIPA) 155.7, 217.4, 358.6. Imports from Taiwan \$201bn in 2025 ($+85$bn); Mexico's computer-equipment exports \$85bn (95\% to the US, 82\% HS 8471.50 assembled servers). Onshoring: TSMC Arizona Blackwell wafers in volume since October 2025 but packaged in Taiwan; Wistron Fort Worth opened July 2026 (GB300 boards); Foxconn Houston ramping; the ``\$500bn of US-made AI infrastructure'' is partner output, not Nvidia capex.
  \item \textbf{Implication for the node.} Under the border rule 80\% of a data-center IT dollar is imported (Fed). Under US+Nvidia, the accounting is a single identity for the node's 2025 IT purchases (about \$200bn: the server share of B935RC plus the networking increment): Nvidia US-installed \$100bn (about 52\% of data-center revenue, below the 69\% customer-HQ share because revenue is recognised at shipment to the ODMs one to two quarters before installation and part of Taiwan-HQ sales is installed abroad), of which 29 is Nvidia's imported cost of revenue, 4 non-US R\&D and 67 US value added; other AI-server content \$40bn and conventional servers \$60bn, each 20\% domestic. Domestic content of a data-center IT dollar: 60\% gross and about 40\% in value added under US+Nvidia, 20\% and about 15\% under the border rule, 100\% in a pure closed economy. The imported part enters as an import row at an exogenous world price financed by a proportional export (balanced trade), not as the fixed-supply device used for fuels. The treatment is Nvidia-only: Broadcom, AMD and Marvell AI silicon, and Google TPU and AWS Trainium (own-account designs that enter US investment at cost, with no market markup) are left as they are. Consistency check: 334's gross operating surplus in the summary Use tables is flat at \$94.5bn (2019), 96.4, 99.3, 99.3, 97.5 (2023) while Nvidia's operating income went 2.8 $\to$ 4.5 $\to$ 10.0 $\to$ 4.4 $\to$ 33.0bn over the same fiscal years, so Nvidia's profits are not in 334's value added either (through the 2023 vintage); the addition to 334 is therefore a gross addition, to be re-checked on the 2024 tables (FY2025 operating income 81bn).
\end{itemize}

\section{Hard steps and open items}\label{sec:hard}
\begin{enumerate}
  \item \textbf{Level of the node.} No official gross output or capital stock for data centers as a function. The level comes from capacity (FERC 50~GW; LBNL TWh) times engineering cost; SAS 518 revenue and the Digital Economy cloud series are cross-checks. Document the sensitivity of Domar weight and upstreamness to this choice.
  \item \textbf{Row re-weighting.} The 518200 row is payment processing and government; cloud revenue by customer industry (IDC, Synergy) is commercial. Alternative: proportional to each sector's investment in computers and software (FA detail), which is official.
  \item \textbf{Import content of the bundle} (Section~\ref{sec:design}): a model decision, not a data one; it halves the domestic multiplier.
  \item \textbf{Shock and its size}: news about future productivity of which sectors, of what size (Acemoglu 0.66\% to Korinek--Suh), with what probability of realisation. Interacts with the P\&P length.
  \item \textbf{Snapshot steady state}: the node's share of investment goes from 1\% of GDP (2023) to 3\% (2027 projections); the steady state is a year, and the boom is a transition. The energy split of \texttt{data\_notes.tex} has the same issue with renewables.
  \item \textbf{Electricity interaction}: the node's marginal electricity is gas-heavy (Virginia, Ohio, Texas) and part of it will be on-site gas; requires the electricity split into fossil and non-fossil generation to matter. The channel from the boom to gas demand is investment in generation (EI21 steam and EI22 internal combustion engines, gas turbines in Machinery/Electrical) and utilities' capex response, not the intermediate share.
  \item \textbf{To retrieve}: NIPA 5.4.5U data center values 2020--2026Q2 (BEA interactive tables; not on FRED); Census C30 historical Excel for 2014--2018 and a consistent monthly 2025--26 path on the September 2026 vintage (\texttt{privsatime.xlsx}); 2024 Use table and GDP by industry (BEA API key); AIES 2023 operating expenses for 518 (electricity line, if published) and the AIES 2023 capex experimental tables (April 2026, 3--4 digit NAICS); Economic Census 2022 product lines for 518210 (colocation, cloud); Digital Economy 2017--2022 cloud spreadsheet. Retrieved: QCEW 518210 (2024 annual), FRED NIPA series, Nvidia 10-K, Census C30 annual by region.
  \item \textbf{Not available}: TFP or output series for data centers; hyperscaler power costs in 10-Ks; official share of on-site gas in data center energy; US capacity by owner type (Synergy US split is paid).
\end{enumerate}

\end{document}
